\documentclass[10pt,a4paper]{amsart}
\usepackage{amsmath,amssymb,amsthm,graphicx}
\usepackage{verbatim}
\newenvironment{reference}{\comment}{\endcomment}
%\newenvironment{reference}{}{}
\newenvironment{slogan}{\comment}{\endcomment}
\newenvironment{history}{\comment}{\endcomment}

% For commutative diagrams we use Xy-pic
\usepackage[all]{xy}

% We use 2cell for 2-commutative diagrams.
\xyoption{2cell}
% Theorem environments.
%
\theoremstyle{plain}
\newtheorem{theorem}[subsection]{Theorem}
\newtheorem{proposition}[subsection]{Proposition}
\newtheorem{lemma}[subsection]{Lemma}

\theoremstyle{definition}
\newtheorem{definition}[subsection]{Definition}
\newtheorem{example}[subsection]{Example}
\newtheorem{exercise}[subsection]{Exercise}
\newtheorem{situation}[subsection]{Situation}

\theoremstyle{remark}
\newtheorem{remark}[subsection]{Remark}
\newtheorem{remarks}[subsection]{Remarks}

\numberwithin{equation}{subsection}

% Macros
%
\def\lim{\mathop{\mathrm{lim}}\nolimits}
\def\colim{\mathop{\mathrm{colim}}\nolimits}
\def\Spec{\mathop{\mathrm{Spec}}}
\def\Hom{\mathop{\mathrm{Hom}}\nolimits}
\def\Ext{\mathop{\mathrm{Ext}}\nolimits}
\def\SheafHom{\mathop{\mathcal{H}\!\mathit{om}}\nolimits}
\def\SheafExt{\mathop{\mathcal{E}\!\mathit{xt}}\nolimits}
\def\Sch{\mathit{Sch}}
\def\Mor{\mathop{\mathrm{Mor}}\nolimits}
\def\Ob{\mathop{\mathrm{Ob}}\nolimits}
\def\Sh{\mathop{\mathit{Sh}}\nolimits}
\def\NL{\mathop{N\!L}\nolimits}
\def\CH{\mathop{\mathrm{CH}}\nolimits}
\def\proetale{{pro\text{-}\acute{e}tale}}
\def\etale{{\acute{e}tale}}
\def\QCoh{\mathit{QCoh}}
\def\Ker{\mathop{\mathrm{Ker}}}
\def\Im{\mathop{\mathrm{Im}}}
\def\Coker{\mathop{\mathrm{Coker}}}
\def\Coim{\mathop{\mathrm{Coim}}}

% Boxtimes
%
\DeclareMathSymbol{\boxtimes}{\mathbin}{AMSa}{"02}

%
% Macros for moduli stacks/spaces
%
\def\QCohstack{\mathcal{QC}\!\mathit{oh}}
\def\Cohstack{\mathcal{C}\!\mathit{oh}}
\def\Spacesstack{\mathcal{S}\!\mathit{paces}}
\def\Quotfunctor{\mathrm{Quot}}
\def\Hilbfunctor{\mathrm{Hilb}}
\def\Curvesstack{\mathcal{C}\!\mathit{urves}}
\def\Polarizedstack{\mathcal{P}\!\mathit{olarized}}
\def\Complexesstack{\mathcal{C}\!\mathit{omplexes}}
% \Pic is the operator that assigns to X its picard group, usage \Pic(X)
% \Picardstack_{X/B} denotes the Picard stack of X over B
% \Picardfunctor_{X/B} denotes the Picard functor of X over B
\def\Pic{\mathop{\mathrm{Pic}}\nolimits}
\def\Picardstack{\mathcal{P}\!\mathit{ic}}
\def\Picardfunctor{\mathrm{Pic}}
\def\Deformationcategory{\mathcal{D}\!\mathit{ef}}
\usepackage{cleveref}
\crefname{originallemma}{original lemma}{original lemmas}
\Crefname{originallemma}{Original lemma}{Original lemmas}
\makeatletter
\newlabel{algebra-lemma-field-finite-type-over-domain}{{00FY}{source}{}{sourcepointer-1}{}}
\newlabel{algebra-lemma-field-finite-type-over-domain@cref}{{[originallemma][0][]00FY}{[1][0][]source}{}{}{}}
\makeatother
\crefname{originallemma}{original lemma}{original lemmas}
\Crefname{originallemma}{Original lemma}{Original lemmas}
\makeatletter
\newlabel{algebra-lemma-finite-over-complete-ring}{{031D}{source}{}{sourcepointer-2}{}}
\newlabel{algebra-lemma-finite-over-complete-ring@cref}{{[originallemma][0][]031D}{[1][0][]source}{}{}{}}
\makeatother
\crefname{originallemma}{original lemma}{original lemmas}
\Crefname{originallemma}{Original lemma}{Original lemmas}
\makeatletter
\newlabel{algebra-lemma-integral-overring-surjective}{{00GQ}{source}{}{sourcepointer-3}{}}
\newlabel{algebra-lemma-integral-overring-surjective@cref}{{[originallemma][0][]00GQ}{[1][0][]source}{}{}{}}
\makeatother
\crefname{originalreference}{original reference}{original references}
\Crefname{originalreference}{Original reference}{Original references}
\makeatletter
\newlabel{lemma-finite-type-finite-type-red}{{0GC5}{source}{}{sourcepointer-4}{}}
\newlabel{lemma-finite-type-finite-type-red@cref}{{[originalreference][0][]0GC5}{[1][0][]source}{}{}{}}
\makeatother
\crefname{originalreference}{original reference}{original references}
\Crefname{originalreference}{Original reference}{Original references}
\makeatletter
\newlabel{lemma-rig-closed-point-relative-residue-field}{{0GGA}{source}{}{sourcepointer-5}{}}
\newlabel{lemma-rig-closed-point-relative-residue-field@cref}{{[originalreference][0][]0GGA}{[1][0][]source}{}{}{}}
\makeatother
\title{Rig-closed points and relative finiteness for adic maps}
\author{The Stacks Project Authors}
\date{stacks-project:restricted.tex:local-snapshot-20261012; GNU FDL 1.2 or later}
\begin{document}
\maketitle
\noindent\textit{Editorial note.} All delivered mathematical text is the author's own, copied from the pinned Stacks Project chapter file restricted.tex; no statement or proof step was rewritten, reordered or removed.\par\medskip
\noindent\textit{Editorial note.} The counted proof is the author's complete argument: the proof proves (1) implies (3) through the finite-length argument and the complete-ring finiteness criterion, (3) implies (2) through integrality and spectral surjectivity, and (2) implies (1) through the residue-field lemma plus the Nullstellensatz\par\medskip
\noindent\textit{Editorial note.} The label lemma-rig-closed-point-relative names the counted statement, and the proof environment immediately following it in the source is delivered verbatim as its proof.\par\medskip
\noindent\textit{Editorial note.} The counted statement is carried by the author's own lemma environment in the chapter source.\par\medskip
\noindent\textit{Allowed standard prerequisites (editorial).}\begin{itemize}
\item rig-closed primes of Noetherian adic rings and the description of B/q as a Noetherian local ring of dimension 1 with its adic topology
\item adicity of the composite A -> B/q when varphi is adic
\item finite length of B/(q + varphi(a)) and the finiteness criterion for complete local rings (Algebra, Lemma algebra-lemma-finite-over-complete-ring)
\item integrality implies surjectivity on spectra (Algebra, Lemma algebra-lemma-integral-overring-surjective)
\item the residue field of a rig-closed point is finitely generated as an algebra over the residue field of the image (lemma-rig-closed-point-relative-residue-field of the same chapter, with lemma-finite-type-finite-type-red)
\item the Hilbert Nullstellensatz: a field finitely generated as an algebra over a subfield is a finite extension (Algebra, Lemma algebra-lemma-field-finite-type-over-domain)
\end{itemize}\medskip
\ifdefined\hypertarget\hypertarget{sourcepointer-1}{}\fi
\noindent\textit{Source reference.} Original lemma 00FY: sources/stacks-chow/tags.txt, source lines 588--588. Prerequisite pointer invoked at proof line(s) 3375; it names a standard fact of a sibling chapter that the counted argument uses, and every step of the counted proof built on it is retained verbatim in the delivered body.\par\medskip
\ifdefined\hypertarget\hypertarget{sourcepointer-2}{}\fi
\noindent\textit{Source reference.} Original lemma 031D: sources/stacks-chow/tags.txt, source lines 3980--3980. Prerequisite pointer invoked at proof line(s) 3360; it names a standard fact of a sibling chapter that the counted argument uses, and every step of the counted proof built on it is retained verbatim in the delivered body.\par\medskip
\ifdefined\hypertarget\hypertarget{sourcepointer-3}{}\fi
\noindent\textit{Source reference.} Original lemma 00GQ: sources/stacks-chow/tags.txt, source lines 616--616. Prerequisite pointer invoked at proof line(s) 3366; it names a standard fact of a sibling chapter that the counted argument uses, and every step of the counted proof built on it is retained verbatim in the delivered body.\par\medskip
\ifdefined\hypertarget\hypertarget{sourcepointer-4}{}\fi
\noindent\textit{Source reference.} Original reference 0GC5: sources/stacks-restricted/pinned-original-source.tex, source lines 2738--2738. Prerequisite pointer invoked at proof line(s) 3372; it names a standard fact of the same chapter that the counted argument uses, and every step of the counted proof built on it is retained verbatim in the delivered body.\par\medskip
\ifdefined\hypertarget\hypertarget{sourcepointer-5}{}\fi
\noindent\textit{Source reference.} Original reference 0GGA: sources/stacks-restricted/pinned-original-source.tex, source lines 3314--3314. Prerequisite pointer invoked at proof line(s) 3371; it names a standard fact of the same chapter that the counted argument uses, and every step of the counted proof built on it is retained verbatim in the delivered body.\par\medskip
\begin{lemma}
\label{lemma-rig-closed-point-relative}
Let $\varphi : A \to B$ be an arrow of $\textit{WAdm}^{Noeth}$
which is adic and topologically of finite type.
Let $\mathfrak q \subset B$ be rig-closed.
Let $\mathfrak p = \varphi^{-1}(\mathfrak q) \subset A$.
Let $\mathfrak a \subset A$ be the ideal of topologically nilpotent
elements.
The following are equivalent
\begin{enumerate}
\item the residue field $\kappa$ of $B/\mathfrak q$ is finite
over $A/\mathfrak a$,
\item $\mathfrak p \subset A$ is rig-closed,
\item $A/\mathfrak p \subset B/\mathfrak q$ is a finite extension
of rings.
\end{enumerate}
\end{lemma}

\begin{proof}
Assume (1). Recall that $B/\mathfrak q$ is a Noetherian local ring
of dimension $1$ whose topology is the adic topology coming
from the maximal ideal. Since $\varphi$ is adic, we see that
$A \to B/\mathfrak q$ is adic. Hence $\varphi(\mathfrak a)$
is a nonzero ideal in $B/\mathfrak q$. Hence
$B/\mathfrak q + \varphi(\mathfrak a)$
has finite length. Hence $B/\mathfrak q + \varphi(\mathfrak a)$
is finite as an $A/\mathfrak a$-module by our assumption.
Thus $B/\mathfrak q$ is finite over $A$ by
Algebra, Lemma \ref{algebra-lemma-finite-over-complete-ring}.
Thus (3) holds.

\medskip\noindent
Assume (3). Then $\Spec(B/\mathfrak q) \to \Spec(A/\mathfrak p)$
is surjective by
Algebra, Lemma \ref{algebra-lemma-integral-overring-surjective}.
This implies (2).

\medskip\noindent
Assume (2). Denote $\kappa'$ the residue field of $A/\mathfrak p$.
By Lemma \ref{lemma-rig-closed-point-relative-residue-field}
(and Lemma \ref{lemma-finite-type-finite-type-red})
the extension $\kappa/\kappa'$ is finitely generated as an algebra.
By the Hilbert Nullstellensatz (Algebra, Lemma
\ref{algebra-lemma-field-finite-type-over-domain})
we see that $\kappa/\kappa'$ is a finite extension.
Hence we see that (1) holds.
\end{proof}

\end{document}
